\section{Datos y metodología}

El desarrollo se organizó comprendiendo el problema de negocio, comprendiendo y preparando los datos, modelando y evaluando. La fase de despliegue se dejó fuera del alcance, dado que el objetivo es comparar alternativas sobre una base reproducible.

\subsection{Fuente y relevamiento de datos}
El conjunto de datos utilizado es \textit{BankSim}, disponible públicamente en Kaggle \cite{dataset}. Fue desarrollado por López-Rojas y Axelsson \cite{inproceedings} como un simulador de pagos bancarios que permite generar transacciones sintéticas para investigaciones sobre detección de fraude.

El archivo analizado contiene 594.643 transacciones y diez columnas, con operaciones distribuidas entre los días 0 y 179 de la simulación. La variable objetivo es \texttt{fraud}, que indica si una transacción fue etiquetada como fraudulenta.

\begin{table}[!htbp]
\centering
\small
\caption{Variables originales del conjunto BankSim}
\label{tab:variables}
\begin{tabular}{@{}p{2.65cm}p{6.0cm}p{3.3cm}@{}}
\toprule
\textbf{Variable} & \textbf{Descripción} & \textbf{Tipo} \\
\midrule
\texttt{step} & Día de la simulación & Temporal \\
\texttt{customer} & Identificador del cliente & Identificador \\
\texttt{age} & Rango de edad del cliente & Categórica \\
\texttt{gender} & Género o tipo de cliente & Categórica \\
\texttt{zipcodeOri} & Código postal del cliente & Categórica \\
\texttt{merchant} & Identificador del comercio & Identificador \\
\texttt{zipMerchant} & Código postal del comercio & Categórica \\
\texttt{category} & Categoría de compra o comercio & Categórica \\
\texttt{amount} & Importe de la transacción & Numérica \\
\texttt{fraud} & Etiqueta de fraude (0 o 1) & Objetivo binario \\
\bottomrule
\end{tabular}
\par\smallskip
\end{table}

\subsection{Exploración inicial y calidad}
El análisis de datos exploratorio (EDA) comenzó con la verificación de los tipos de variables, las dimensiones del conjunto y la presencia de valores faltantes. No se observaron valores nulos en las columnas conservadas para el modelado. Los textos se normalizaron eliminando las comillas simples que delimitaban sus valores.

Se eliminaron \texttt{age} y \texttt{gender}, dado que los autores de BankSim aclaran que el comportamiento de consumo de los agentes no se programó en función de estas características \cite{inproceedings}. También se descartaron \texttt{zipcodeOri} y \texttt{zipMerchant}, que en el anteproyecto se identificaron como variables sin variación. Tras esta etapa quedaron seis columnas originales, incluida la variable objetivo.

Como aproximación inicial a la relación con el fraude, se obtuvo una correlación de Pearson de 0,490 entre \texttt{amount} y \texttt{fraud}. Para las categorías, analizadas como indicadores binarios, se observaron asociaciones positivas en \texttt{es\_sportsandtoys} (0,364), \texttt{es\_travel} (0,250) y \texttt{es\_leisure} (0,248), y negativa en \texttt{es\_transportation} ($-0,263$). Estas relaciones son descriptivas del dataset y no implican causalidad. Asimismo, no se encontraron valores atípicos que pertenezcan a errores, por lo que se decidio conservar la mayor cantidad de información.

\subsection{Ingeniería de variables}
Siguiendo los lineamientos del manual de detección de fraude \cite{leborgne2022fraud}, se construyeron dieciocho variables adicionales para incorporar información histórica en ventanas de 1, 7 y 30 días. El procedimiento consideró tres niveles de agregación:

\begin{table}[!htbp]
\centering
\small
\caption{Variables generadas mediante \textit{feature engineering}}
\label{tab:features}
\begin{tabular}{@{}p{2.65cm}p{6.55cm}p{6cm}@{}}
\toprule
\textbf{Entidad} & \textbf{Variables} & \textbf{Información representada} \\
\midrule
Cliente & \texttt{tx\_count\_last\_nd}; \texttt{avg\_amount\_last\_nd} & Cantidad y monto promedio de transacciones recientes. \\
Comercio & \texttt{merchant\_tx\_delayed\_nd}; \texttt{merchant\_risk\_nd} & Actividad y proporción de fraudes en antecedentes disponibles. \\
Categoría & \texttt{category\_tx\_delayed\_nd}; \texttt{category\_risk\_nd} & Actividad y proporción de fraude de la categoría. \\
\bottomrule
\end{tabular}
\par\smallskip
{\footnotesize En los nombres, \texttt{n} toma los valores 1, 7 y 30.}
\end{table}

Para calcular el riesgo histórico de un comercio o categoría se utilizó la relación entre fraudes observados y transacciones registradas en una ventana anterior:
\begin{equation}
\label{eq:riesgo}
\operatorname{Riesgo}_{n} = \frac{\text{Fraudes históricos en }n\text{ días}}{\text{Transacciones históricas en }n\text{ días}}.
\end{equation}

Las ventanas de riesgo excluyen el día que se está evaluando y los siete días previos: para una transacción del día $t$, los antecedentes utilizados llegan hasta $t-8$. En ausencia de transacciones en la ventana, el indicador de riesgo se definió como cero. Las ventanas de cliente, en cambio, incluyen la operación actual; esta particularidad debe tenerse en cuenta si el procedimiento se adapta a un entorno de detección en tiempo real con marcas temporales más precisas.

La idea de construir indicadores de importe anómalo respecto del comportamiento individual se consideró, pero no se incorporó a los experimentos finales. Después de la ingeniería de variables, el conjunto contó con 24 columnas antes de aplicar codificación categórica.

\subsection{División temporal y filtrado de clientes}
La separación se realizó por días, evitando mezclar observaciones anteriores y posteriores entre entrenamiento y evaluación. Se dejaron dos intervalos de siete días que representan el retraso de disponibilidad de las etiquetas.

\begin{table}[!htbp]
\centering
\small
\caption{Partición temporal y efecto del filtrado}
\label{tab:splits}
\begin{tabular}{@{}lrrrr@{}}
\toprule
\textbf{Período} & \textbf{Días} & \textbf{Transacciones} & \textbf{Fraudes} & \textbf{Tras filtrado} \\
\midrule
Entrenamiento & 0--119 & 374.914 & 4.800 & 374.914 \\
Espera & 120--126 & 24.986 & 280 & No evaluado \\
Validación & 127--149 & 83.437 & 920 & 60.655 \\
Espera & 150--156 & --- & --- & No evaluado \\
Prueba & 157--179 & 85.380 & 920 & 58.305 \\
\bottomrule
\end{tabular}
\par\smallskip
{\footnotesize Las columnas de transacciones y fraudes corresponden a los periodos sin filtrar.}
\end{table}

Además de respetar el orden temporal, se eliminaron de validación y prueba las operaciones de clientes cuyo fraude ya debía ser conocido al momento de emitir la alerta. El criterio supone que \textit{una entidad financiera dejaría de investigar una cuenta que ya fue identificada y bloqueada}. Se incorporaron progresivamente al conjunto de clientes conocidos los fraudes ocurridos hasta el día $t-8$.

El entrenamiento incluyó 4.111 clientes diferentes, de los cuales 1.166 presentaron al menos un fraude. El filtrado redujo las operaciones de validación de 83.437 a 60.655 y los fraudes de 920 a 160. En prueba, las operaciones se redujeron de 85.380 a 58.305 y los fraudes de 920 a 113. En los conjuntos filtrados se identificaron 142 clientes distintos con fraude en validación y 104 en prueba.

La proporción de operaciones fraudulentas resultante fue de 1,280\% en entrenamiento, 0,264\% en validación y 0,194\% en prueba. Estas prevalencias reflejan, entre otros factores, el filtrado aplicado; por lo tanto, no deben compararse como si fueran muestras aleatorias equivalentes.

\subsection{Preprocesamiento para el modelado}
Se excluyeron \texttt{fraud}, \texttt{customer} y \texttt{step} de las variables predictoras: la primera es la etiqueta, mientras que las otras dos se conservaron para organizar el análisis temporal y calcular las métricas por cliente. Se utilizó un \texttt{ColumnTransformer} con \texttt{StandardScaler} para las variables numéricas y \texttt{OneHotEncoder} para \texttt{merchant} y \texttt{category}. El transformador se ajustó exclusivamente con el conjunto de entrenamiento y luego se aplicó a validación y prueba, con tratamiento de categorías no observadas durante el ajuste.